\begin{textsample}{POS Dim 2 – ac – Score 34.00 – ac/ac\_em\_35.txt}  \label{ex:f2_pos_162}
18. \textbf{Legislação} A \textbf{legislação} e os documentos norteadores utilizados na elaboração e na organização do documento curricular que contemplem a Formação \textbf{Técnica} e Profissional atendem às \textbf{prerrogativas} da Lei nº 13.415, de 16 de \textbf{fevereiro} de 2017, alterando a Lei n{\EmojiFont °} 9.394, de 20 de dezembro de 1996, que, por sua vez, estabelece as \textbf{diretrizes} e bases da educação nacional e dá outras \textbf{providências}.
Outras normas e \textbf{diretrizes} orientaram todo esse processo de mudanças, tais como as \textbf{Diretrizes} Curriculares Nacionais para o Ensino Médio, \textbf{publicadas} na Resolução MEC/CNE/CEB nº 3, de 21 de \textbf{novembro} de 2018 ; Resolução CNE/CEB nº 1, de 5 de dezembro de 2014, que \textbf{atualiza} e define novos critérios para a composição do \textbf{Catálogo} Nacional de \textbf{Cursos} \textbf{Técnicos}, disciplinando e orientando os sistemas de ensino e as \textbf{instituições} públicas e privadas de Educação Profissional e Tecnológica quanto à \textbf{oferta} de \textbf{cursos} \textbf{técnicos} de nível médio ; a Lei n{\EmojiFont °} 13.005 de 25 de junho de 2014 – que aprova o Plano Nacional de Educação - PNE e dá outras \textbf{providências} ; a Lei nº 11.741, de 16 de \textbf{julho} de 2008, que altera dispositivos da redação original da LDB, para redimensionar, institucionalizar e integrar as ações da educação profissional \textbf{técnica} de nível médio, da educação de jovens e adultos e da educação profissional e tecnológica ; a Resolução CNE/CP nº 1, de 05 de \textbf{janeiro} de 2021 – responsável por definir as \textbf{Diretrizes} Curriculares Nacionais Gerais para a Educação Profissional e Tecnológica ; a Resolução CEE/AC Nº 177/2013, que dispõe sobre \textbf{diretrizes} gerais para a Educação Profissional \textbf{Técnica} de Nível Médio no âmbito do sistema de ensino do estado do Acre ; e a Resolução Nº 01/2021 de 25 de maio de 2021, que \textbf{institui} \textbf{Diretrizes} Operacionais para a Educação de Jovens e Adultos nos aspectos relativos ao seu alinhamento à Política Nacional Comum Curricular ( BNCC ), e Educação de Jovens e Adultos a Distância. 19.
Organização da \textbf{oferta} A educação profissional e tecnológica é uma modalidade educacional \textbf{prevista} na Lei de \textbf{Diretrizes} e Bases da Educação Nacional ( LDB ) e tem por finalidade preparar o cidadão para o exercício de profissões, para que este possa ser inserido no mundo do trabalho e na vida em sociedade.
Dessa forma, traduz o processo formativo em \textbf{cursos} de \textbf{qualificação}, \textbf{habilitação} \textbf{técnica} e tecnológica e de pós-graduação, organizados de forma a propiciar o aproveitamento contínuo e articulado dos estudos, conforme Portal Brasileiro de Dados Abertos do Ministério da Educação e Cultura ( MEC, 2019 ).
Logo, o estudante que optar pela Formação \textbf{Técnica} e Profissional ( FTP ) poderá fazer um \textbf{curso} \textbf{técnico} e/ou \textbf{cursos} de \textbf{Qualificação} Profissional articulados entre si, sendo mais uma alternativa para os estudantes acreanos.
Coerentes com essa perspectiva, as \textbf{Diretrizes} Curriculares Nacionais para o Ensino Médio ( DCNEM ), \textbf{atualizadas} pelo Conselho Nacional de Educação - CNE em \textbf{novembro} de 2018, indicam que os \textbf{currículos} dessa etapa de ensino devem ser compostos por : >"Formação Geral Básica : Conjunto de competências e habilidades das Áreas de Conhecimento ( Linguagens e suas Tecnologias, Matemática e suas Tecnologias, Ciências da Natureza e suas Tecnologias, Ciências Humanas e Sociais Aplicadas ) \textbf{previstas} na etapa do Ensino Médio da Base Nacional Comum Curricular - BNCC, que aprofundam e consolidam as aprendizagens essenciais do Ensino Fundamental, a compreensão de problemas complexos e a reflexão sobre soluções para eles, com \textbf{carga} \textbf{horária} total máxima de 1.800 horas ” ; \textbf{Itinerários} Formativos : Conjunto de situações e atividades educativas que os estudantes podem escolher conforme seu interesse, para aprofundar e ampliar aprendizagens em uma ou mais Áreas de Conhecimento e/ou na Formação \textbf{Técnica} e Profissional, com \textbf{carga} \textbf{horária} total mínima de 1.200 horas.
Diante do exposto, os estudantes matriculados no Ensino Médio regular da rede de ensino do estado do Acre terão a possibilidade de escolher no Itinerário Formativo as seguintes rotas : I - Rota de Linguagens e suas Tecnologias ; II - Rota de Matemática e suas Tecnologias ; III - Rota de Ciências da Natureza e suas Tecnologias ; IV - Rota de Ciências Humanas e Sociais Aplicadas ; V - Formação \textbf{Técnica} e Profissional.
As Rotas de Aprofundamento relacionadas às áreas de conhecimento possibilitam aos estudantes expandir os aprendizados promovidos pela Formação Geral Básica e prosseguir no nível superior.
A Formação \textbf{Técnica} e Profissional ( FTP ), com abordagem do trabalho como princípio educativo, possibilita a \textbf{qualificação} nas atuais demandas profissionais do mercado de trabalho ( MEC, 2017 ), uma vez que a escola se preocupa com a identificação e definição de saberes e competências profissionais, articuladas às dez competências da BNCC, que devem compor o perfil do aluno egresso do Ensino Médio, garantindo que estejam aptos para as demandas do mercado de trabalho regional e, consequentemente, para as novas exigências ocupacionais geradas pelas próprias transformações no mundo do trabalho, mediante a aproximação da escola e dos estudantes com os diversos setores da sociedade para promover parcerias e vivências práticas.
Dessa forma, espera -se que o perfil egresso desses estudantes que realizaram os \textbf{cursos} da Formação \textbf{Técnica} e Profissional tenham competências, habilidades e atitudes que os prepare para o exercício profissional na área de formação, com visão crítica, humanística pautado pela ética e responsabilidade social.
A \textbf{instituição} parceira que \textbf{oferta} Educação Profissional \textbf{Técnica} de Nível Médio, no âmbito do estado do Acre, deverá realizar, periodicamente, pesquisa de egressos e avaliações dos seus serviços, visando garantir a qualidade da \textbf{oferta}.
A Formação \textbf{Técnica} e Profissional ( FTP ) será disponibilizada para os vinte e dois \textbf{municípios} do Acre por meio de parcerias com \textbf{instituições} credenciadas que já atuam com formação profissional no estado, em \textbf{conformidade} com o termo de cooperação celebrado entre elas.
O portfólio de \textbf{oferta} contempla \textbf{cursos} de \textbf{Qualificação} Profissional e \textbf{cursos} com \textbf{habilitação} \textbf{técnica} e poderá ser modificado no decorrer dos anos, conforme os arranjos produtivos locais e as mudanças do mundo do trabalho.
A organização do portfólio de \textbf{oferta} foi baseada em estudos realizados pelas \textbf{instituições} parceiras responsáveis pela \textbf{oferta} de Educação Profissional no estado, considerando o perfil do mercado de trabalho, as transformações e mudanças que estão ocorrendo no mundo do trabalho, pensando nas áreas e profissões que vão demandar necessidades profissionais.
Nesse sentido, as \textbf{instituições} parceiras responsáveis pela Formação \textbf{Técnica} e Profissional ( FTP ) devem propiciar uma formação humana integral para além da \textbf{oferta} de um conjunto de conhecimentos científicos, \textbf{técnicos} e tecnológicos produzidos e acumulados pela humanidade, sendo capaz de promover o pensamento crítico como possibilidade para compreender o mundo ( da vida e do trabalho ), as ideias, os problemas, as dificuldades e as potencialidades da sociedade, contribuindo para a produção de novos saberes e práticas, voltados para os interesses sociais e coletivos.
Desse modo, a Secretaria de Estado de Educação do Estado do Acre deve acompanhar e garantir a articulação entre a Formação \textbf{Técnica} e Profissional ( FTP ) e a Base Nacional Comum Curricular, oferecida aos estudantes do Ensino Médio, monitorando e avaliando, portanto, além do trabalho realizado pelas escolas de origem do estudante, também as \textbf{instituições} parceiras, além de observar se os Projetos Pedagógicos de \textbf{Cursos} ( PPCs ) foram elaborados de forma coletiva e democrática, na perspectiva de formação para a cidadania e para o mundo do trabalho, garantia da qualidade do ensino, bem como se atende às orientações da \textbf{legislação} \textbf{vigente}, tanto em âmbito nacional como estadual.

% matched lemmas: atualizar, carga, catálogo, conformidade, currículo, curso, diretriz, fevereiro, habilitação, horário, instituir, instituição, itinerário, janeiro, julho, legislação, município, novembro, oferta, prerrogativa, prever, providência, publicar, qualificação, técnico, vigente
\end{textsample}
